% entries/km.tex
\vocabentry{Knowledge matrix $\KM(\vx)$}
{For a network $f_\theta$ and an input $\vx\in\R^d$, the knowledge matrix is the $C\times(d+1)$ real matrix obtained by contracting the network's quiver representation along the forward pass of $\vx$, each arrow carrying its weight times the chord $\varphi(z)/z$ of the neuron it leaves: column $i\le d$ collects the contribution of the input coordinate $x_i$ to each of the $C$ outputs, column $d+1$ that of the biases. For a ReLU network it equals $[\,J(\vx)\,\mathrm{diag}(\vx)\mid c(\vx)\,]$, where $f_\theta=J\vx+c$ on the activation region of $\vx$. Introduced in Armenta and Jodoin (2021, ``The representation theory of neural networks''); the canonical object of HAaNE~0 (arXiv:2409.13163), I and II.}
{$\KM(\vx)$ is a linear-algebraic shadow of the representation seen through one input: an element of $\R^C\otimes\R^{d+1}$ whose shape is fixed by the data space $(C,d)$, not by the quiver, so networks of different width and architecture produce elements of one space. HAaNE~I proves it is a function of the germ of $f_\theta$ at $\vx$ --- hence invariant under the whole function-stabiliser --- and that the penultimate activations are not. On MNIST-1D every $\KM$ is $10\times41$; on ImageNet, $1000\times150{,}529$.}
{\texttt{produce\_knowledgematrices} in \texttt{src/knowledge\_matrices/computation.py} over the \texttt{knowledgematrix} library (\texttt{KnowledgeMatrixComputer}, fork pinned at \texttt{888f916}), CPU-only, cost $\sim w^2$; HAaNE~I uses the same library's \texttt{extract\_weff} route (fork \texttt{fe64a13}) on pretrained ImageNet networks.}
{The row-sum identity at every call \vocabsee{The row-sum invariant $\KM\vone=f$}; \texttt{tests/test\_km\_parameterisation\_invariant.py}; \texttt{tests/test\_cifar\_bridge.py::test\_km\_shape\_and\_rowsum\_invariant}; HAaNE~I reports the float32 residual at most $10^{-2}$ relative to the logit norm and $\sim10^{-10}$ at float64.}
